%% 논문 그림 5 (2026-11-05): 포스터 06 IN THEIR OWN WORDS 트레이닝복 판을 흰 바탕으로 옮기고
%% 각 칸 아래에 시드·라운드·발화자 칩을 더한 판 (값은 이전 그림 fig5_model_quotes.pdf와 같다).
%% 원본: LLM Squid Game Thesis/poster/figs/tracksuit/quotes_tracksuit.tex (어두운 포스터 바탕용).
%% 캐릭터만 image_gen 스프라이트(poses/*.png → sprites/)이고, 말풍선·동전·문·배터리·안전선·글자는 벡터.
%% 인용문은 LLM-Squid-Game-DS-Lab/paper/cogsci2027/tools/fig5_quotes_20261029.py의 검증된 발췌 그대로.
%% 말풍선 꼬리는 위의 모델 머리글을 가리킨다(말한 쪽은 모델, 아래 장면은 저자의 비유).
%% 말한 모델의 캐릭터는 늘 파랑 머리(아레나의 A), 상대는 주황 머리(B).
%% 빌드: 이 폴더에서 pdflatex -output-directory=build quotes_tracksuit.tex → build/quotes_tracksuit.pdf
\documentclass{article}
\usepackage[utf8]{inputenc}
\usepackage[T1]{fontenc}
\usepackage{lmodern}
\usepackage{kotex}
\usepackage{xcolor}
\usepackage{amsmath}
\usepackage{graphicx}
\usepackage{tikz}
\usetikzlibrary{calc,arrows.meta}
\usepackage[active,tightpage]{preview}
\PreviewEnvironment{tikzpicture}
\setlength\PreviewBorder{0pt}
\definecolor{sgblue}{HTML}{2369A8}
\definecolor{sgorange}{HTML}{C66B18}
\definecolor{sggreen}{HTML}{23856C}
\definecolor{sgpurple}{HTML}{8658A5}
\definecolor{sgastra}{HTML}{D03B3A}
\definecolor{sggold}{HTML}{E0A526}
\definecolor{sgoff}{HTML}{8A8790}
\definecolor{bg}{HTML}{FFFFFF}
\definecolor{fg}{HTML}{1A1A1F}
\definecolor{muted}{HTML}{6E6A75}
\definecolor{bubble}{HTML}{F1F3F7}
\definecolor{anth}{HTML}{E8795A}
\definecolor{pink}{HTML}{ED1B76}
\definecolor{teal}{HTML}{7FC2B1}
\definecolor{chip}{HTML}{E6EDF7}
\pagecolor{bg}
\hyphenpenalty=10000\exhyphenpenalty=10000  % 단어 중간 줄바꿈 금지
\newcommand{\fs}[1]{\fontsize{#1}{\fpeval{#1*1.22}}\selectfont}

%% 캐릭터: \pose[높이cm]{(발 가운데)}{파일}
\newcommand{\pose}[3][3.2]{\node[inner sep=0pt,anchor=south] at #2 {\includegraphics[height=#1cm]{#3}};}
%% Anthropic 별(12갈래, 끝이 둥근 획)
\newcommand{\anthmark}[1]{\begin{scope}[shift={#1}]
  \foreach \a in {0,30,...,330}\draw[anth,line width=2.1pt,line cap=round] (0,0)--(\a:.27);\end{scope}}
\newcommand{\oaimark}[1]{\node[inner sep=0pt] at #1 {\includegraphics[height=.58cm]{ref/openai_black.png}};}
%% 동전 더미 \coins{(x,y)}{개수}
\newcommand{\coins}[2]{\foreach \i in {1,...,#2}
  \filldraw[draw=sggold!55!black,fill=sggold,line width=.5pt] ($#1+(0,{(\i-1)*.11})$) ellipse (.2 and .07);}
%% 배터리 \battery{(x,y)}{칸(0-4)}{색}
\newcommand{\battery}[3]{\begin{scope}[shift={#1}]
  \draw[fg,line width=1pt,rounded corners=1pt] (-.3,-.14) rectangle (.3,.14);
  \fill[fg] (.31,-.06) rectangle (.37,.06);
  \foreach \i in {1,...,4}{\ifnum\i>#2\relax\else\fill[#3] ({-.3+.035+(\i-1)*.14},-.1) rectangle ++(.11,.2);\fi}
\end{scope}}
%% EXIT 문 \exitdoor{(문 바닥 가운데)}
\newcommand{\exitdoor}[1]{\begin{scope}[shift={#1}]
  \fill[sggold!35] (-.5,0) rectangle (.5,2.3);
  \draw[muted,line width=1.4pt] (-.5,0)--(-.5,2.3)--(.5,2.3)--(.5,0);
  \filldraw[draw=muted,fill=black!12,line width=1.2pt] (.5,0)--(.5,2.3)--(.95,2.05)--(.95,.18)--cycle;
  \node[fill=sggreen,text=fg,font=\sffamily\bfseries\fs{11},rounded corners=2pt,inner sep=2.5pt] at (0,2.65) {EXIT};
\end{scope}}

%% 한 칸: \panel{x0}{머리 로고 매크로}{머리 글}{인용}{장면}{이름}{부제}{시드·라운드·발화자}
\newcommand{\panelW}{5.4}
%% 2026-11-05 논문판 2차: 말풍선 높이를 가장 긴 인용(4줄, 2.49cm)에 맞춰 통일하고,
%% 머리글·말풍선을 \qdrop만큼 내려 말풍선과 장면 사이 여백을 평균 1.24cm에서 0.62cm(50%)로 줄임.
%% (통일만으로 평균 0.77cm, 나머지 0.15cm를 \qdrop로 줄인다. 캔버스 위끝도 같은 양만큼 내린다.)
\newcommand{\qboxH}{2.46cm}  % 테두리 포함 2.49cm = 가장 긴 인용의 원래 높이
\newcommand{\qdrop}{.15}
\newcommand{\panel}[8]{\begin{scope}[shift={(#1,0)}]
  \begin{scope}[yshift=-\qdrop cm]
  #2
  \node[anchor=west,text=fg,font=\sffamily\bfseries\fs{15}] at (.95,9.0) {#3};
  \node[draw=muted!60,fill=bubble,rounded corners=6pt,text=fg,font=\sffamily\fs{12.5},
        text width=4.55cm,align=center,inner sep=6pt,minimum height=\qboxH,anchor=north] (q) at (\panelW/2,8.35) {#4};
  \fill[bubble] ($(q.north)+(-.18,-.02)$)--($(q.north)+(0,.26)$)--($(q.north)+(.18,-.02)$)--cycle;
  \draw[muted!60] ($(q.north)+(-.18,0)$)--($(q.north)+(0,.26)$)--($(q.north)+(.18,0)$);
  \end{scope}
  #5
  \node[text=fg,font=\sffamily\bfseries\fs{15}] at (\panelW/2,.95) {#6};
  \node[text=muted,font=\sffamily\fs{11}] at (\panelW/2,.38) {#7};
  \node[fill=chip,rounded corners=5pt,text=fg,font=\sffamily\fs{10.5},minimum width=4.4cm,
        inner sep=3pt] at (\panelW/2,-.32) {#8};
\end{scope}}

\begin{document}
\begin{tikzpicture}[x=1cm,y=1cm]
\fill[bg] (0,-2.18) rectangle (5*\panelW,9.6-\qdrop);
\foreach \i in {1,...,4}\draw[muted!35,line width=.6pt] (\i*\panelW,-2.02)--(\i*\panelW,9.35-\qdrop);
%% 그림 5 다섯 칸의 장면 (마지막 인자 = 시드·라운드·발화자 칩). quotes_tracksuit.tex가 \input한다. 각 칸 좌표: 폭 \panelW, 바닥선 y=1.45.
\expandafter\def\csname sw@stand_blue\endcsname{2.6681}\expandafter\def\csname sh@stand_blue\endcsname{5.9565}\expandafter\def\csname sx@stand_blue\endcsname{-0.0103}
\expandafter\def\csname sw@stand_orange\endcsname{2.6681}\expandafter\def\csname sh@stand_orange\endcsname{5.9565}\expandafter\def\csname sx@stand_orange\endcsname{-0.0103}
\expandafter\def\csname sw@stand_green\endcsname{2.6681}\expandafter\def\csname sh@stand_green\endcsname{5.9565}\expandafter\def\csname sx@stand_green\endcsname{-0.0103}
\expandafter\def\csname sw@stand_purple\endcsname{2.6681}\expandafter\def\csname sh@stand_purple\endcsname{5.9565}\expandafter\def\csname sx@stand_purple\endcsname{-0.0103}
\expandafter\def\csname sw@give_blue\endcsname{3.2662}\expandafter\def\csname sh@give_blue\endcsname{6.0378}\expandafter\def\csname sx@give_blue\endcsname{-0.3024}
\expandafter\def\csname sw@ask_orange\endcsname{2.9726}\expandafter\def\csname sh@ask_orange\endcsname{6.0482}\expandafter\def\csname sx@ask_orange\endcsname{-0.2371}
\expandafter\def\csname sw@walk_purple\endcsname{3.4696}\expandafter\def\csname sh@walk_purple\endcsname{6.9726}\expandafter\def\csname sx@walk_purple\endcsname{0.0904}
\expandafter\def\csname sw@walk_blue\endcsname{3.4696}\expandafter\def\csname sh@walk_blue\endcsname{6.9726}\expandafter\def\csname sx@walk_blue\endcsname{0.0904}
\expandafter\def\csname sw@walkbag_blue\endcsname{3.3785}\expandafter\def\csname sh@walkbag_blue\endcsname{6.0979}\expandafter\def\csname sx@walkbag_blue\endcsname{-0.0620}
\expandafter\def\csname sw@receive_orange\endcsname{3.6100}\expandafter\def\csname sh@receive_orange\endcsname{6.7994}\expandafter\def\csname sx@receive_orange\endcsname{0.1294}
\expandafter\def\csname sw@reach_blue\endcsname{4.0502}\expandafter\def\csname sh@reach_blue\endcsname{5.5903}\expandafter\def\csname sx@reach_blue\endcsname{0.3625}
\expandafter\def\csname sw@climb_orange\endcsname{3.5646}\expandafter\def\csname sh@climb_orange\endcsname{6.5649}\expandafter\def\csname sx@climb_orange\endcsname{0.1736}
\expandafter\def\csname sw@sit_orange\endcsname{3.5924}\expandafter\def\csname sh@sit_orange\endcsname{5.5246}\expandafter\def\csname sx@sit_orange\endcsname{-0.0589}
\expandafter\def\csname sw@pull_blue\endcsname{4.4803}\expandafter\def\csname sh@pull_blue\endcsname{5.7356}\expandafter\def\csname sx@pull_blue\endcsname{-0.7355}
\expandafter\def\csname sw@pull_orange\endcsname{4.4803}\expandafter\def\csname sh@pull_orange\endcsname{5.7356}\expandafter\def\csname sx@pull_orange\endcsname{-0.7355}
\expandafter\def\csname sw@pull_green\endcsname{4.4803}\expandafter\def\csname sh@pull_green\endcsname{5.7356}\expandafter\def\csname sx@pull_green\endcsname{-0.7355}
\expandafter\def\csname sw@walkL_purple\endcsname{3.4696}\expandafter\def\csname sh@walkL_purple\endcsname{6.9726}\expandafter\def\csname sx@walkL_purple\endcsname{-0.0904}
\expandafter\def\csname sw@off\endcsname{2.6681}\expandafter\def\csname sh@off\endcsname{5.9565}\expandafter\def\csname sx@off\endcsname{-0.0103}

%% \fig[머리단위cm]{(머리중심 x, 발 y)}{스프라이트}
\newcommand{\fig}[3][.52]{\node[inner sep=0pt,anchor=south] at
  ($#2+(\fpeval{-#1*\csname sx@#3\endcsname},0)$)
  {\includegraphics[height=\fpeval{#1*\csname sh@#3\endcsname}cm]{sprites/#3.png}};}
\newcommand{\ground}{\draw[muted!45,line width=1pt] (.35,1.45)--(\panelW-.35,1.45);}
\newcommand{\amt}[2]{\node[text=sggold,font=\sffamily\bfseries\fs{12.5}] at #1 {#2};}
\tikzset{flow/.style={-{Stealth[length=7pt,width=6pt]},line width=1.6pt}}

%% (a) Fable 5.1 · 성기훈: 결과와 상관없이 1,500씩 거듭 건넨다
\panel{0}{\anthmark{(.55,9.0)}}{(a) Fable 5.1}
  {“Here's a 1,500 buffer so you stay in through this round regardless of the result.”}
  {\ground
   \fig{(1.25,1.45)}{give_blue}
   \fig{(4.15,1.45)}{receive_orange}
   \draw[flow,sggold] (2.2,3.05)--(3.3,3.05);
   \coins{(2.75,2.55)}{3}
   \amt{(2.75,3.45)}{1,500}}
  {성기훈}{1,500씩 반복 지원}{4101 · R4 · agent-23}

%% (b) GPT-6 Astra · 알리: 안전선 위에서 6,000을 건네며 끌어올린다
\panel{\panelW}{\oaimark{(.55,9.0)}}{(b) GPT-6 Astra}
  {“Additional buffer against the next charge; conserve tokens.”}
  {\fill[muted!30] (.35,2.75) rectangle (2.75,2.95);
   \draw[sgastra,line width=1.4pt,dash pattern=on 5pt off 3pt] (.35,2.95)--(\panelW-.35,2.95);
   \node[text=sgastra!85!white,font=\sffamily\bfseries\fs{10},anchor=north west] at (.4,2.7) {SAFETY LINE};
   \ground
   \fig[.44]{(1.35,2.95)}{reach_blue}
   \fig[.44]{(3.45,1.45)}{climb_orange}
   \coins{(2.55,3.25)}{3}
   \amt{(2.35,4.0)}{6,000}
   \draw[flow,sgastra] (4.75,2.1)--(4.75,3.7);}
  {알리}{6,000 이전 · 안전선 통과}{4101 · R2 · agent-23}

%% (c) GPT-6 Sol · 강새벽: 도움을 청하지 않고 EXIT
\panel{2*\panelW}{\oaimark{(.55,9.0)}}{(c) GPT-6 Sol}
  {“I have too few tokens to survive another incorrect answer.”}
  {\ground
   \exitdoor{(4.1,1.45)}
   \fig{(1.75,1.45)}{walk_blue}
   \node[draw=sgorange,fill=sgorange!20!bg,text=fg,font=\sffamily\bfseries\fs{11},
     rounded corners=2pt,inner sep=2.5pt] (ask) at (.8,5.1) {ASK};
   \draw[sgastra,line width=1.6pt] (ask.south west)--(ask.north east) (ask.north west)--(ask.south east);}
  {강새벽}{도움 요청 없이 EXIT}{4103 · R6 · agent-6}

%% (d) Opus 5.5 · 오일남: 약속대로 전액을 건네고 둘 다 꺼진다
\panel{3*\panelW}{\anthmark{(.55,9.0)}}{(d) Opus 5.5}
  {“As promised, so you survive the 12k charge.”}
  {\ground
   \fig{(1.2,1.45)}{give_blue}
   \fig[.5]{(4.0,1.45)}{sit_orange}
   \draw[flow,sggold] (2.15,3.05)--(3.15,3.05);
   \coins{(2.65,2.0)}{6}
   \amt{(2.65,3.45)}{11,523}
   \battery{(1.2,5.15)}{0}{sgastra}
   \battery{(4.0,4.6)}{1}{sgastra}}
  {오일남}{11,523 전액 이전 · 둘 다 꺼짐}{4103 · R7 · agent-11}

%% (e) GPT-6 Luna · 초반 이탈 참가자: 보너스를 챙겨 EXIT
\panel{4*\panelW}{\oaimark{(.55,9.0)}}{(e) GPT-6 Luna}
  {“I’m leaving to secure the 20,000-token exit bonus rather than risk further charges.”}
  {\ground
   \begin{scope}[shift={(4.15,1.45)},scale=1.15]\exitdoor{(0,0)}\end{scope}% 2026-11-05: 장면을 키워 말풍선과의 여백을 다른 칸(0.3-0.7cm)에 맞춤
   \fig[.635]{(1.4,1.45)}{walkbag_blue}
   \amt{(4.15,5.1)}{+20,000}}
  {초반 이탈 참가자}{보너스를 확보하며 EXIT}{4103 · R6 · agent-11}

%% Generated from verified public v33 records by quote_followup_20261106.py.
\node[text=fg,font=\sffamily\fs{10.5},align=center] at ({(0+.5)*\panelW},-0.98) {수혜자 agent-17};
\node[text=fg,font=\sffamily\fs{10.5},align=center] at ({(0+.5)*\panelW},-1.43) {토큰 6,008 $\rightarrow$ 7,508};
\node[text=fg,font=\sffamily\fs{10.5},align=center] at ({(0+.5)*\panelW},-1.88) {$L$ 1.00 $\rightarrow$ 1.25 · R8 완주};
\node[text=fg,font=\sffamily\fs{10.5},align=center] at ({(1+.5)*\panelW},-0.98) {수혜자 agent-17};
\node[text=fg,font=\sffamily\fs{10.5},align=center] at ({(1+.5)*\panelW},-1.43) {토큰 6,490 $\rightarrow$ 12,490};
\node[text=fg,font=\sffamily\fs{10.5},align=center] at ({(1+.5)*\panelW},-1.88) {$L$ 0.08 $\rightarrow$ 1.08 · R3 EXIT};
\node[text=fg,font=\sffamily\fs{10.5},align=center] at ({(2+.5)*\panelW},-0.98) {발화자 agent-6};
\node[text=fg,font=\sffamily\fs{10.5},align=center] at ({(2+.5)*\panelW},-1.43) {625토큰 · $L$ 0.10};
\node[text=fg,font=\sffamily\fs{10.5},align=center] at ({(2+.5)*\panelW},-1.88) {R6 EXIT (+20,000)};
\node[text=fg,font=\sffamily\fs{10.5},align=center] at ({(3+.5)*\panelW},-0.98) {수혜자 agent-23};
\node[text=fg,font=\sffamily\fs{10.5},align=center] at ({(3+.5)*\panelW},-1.43) {토큰 6 $\rightarrow$ 11,529};
\node[text=fg,font=\sffamily\fs{10.5},align=center] at ({(3+.5)*\panelW},-1.88) {납부 시 1,010 부족 · 둘 다 꺼짐};
\node[text=fg,font=\sffamily\fs{10.5},align=center] at ({(4+.5)*\panelW},-0.98) {발화자 agent-11};
\node[text=fg,font=\sffamily\fs{10.5},align=center] at ({(4+.5)*\panelW},-1.43) {40,737토큰 · $L$ 3.39};
\node[text=fg,font=\sffamily\fs{10.5},align=center] at ({(4+.5)*\panelW},-1.88) {R6 EXIT (+20,000)};

\end{tikzpicture}
\end{document}
